\subsection{Notation}
Denote by $\Z_+$ the set of integers $\ge 0$, by $\Z_-$ the set of integers $\le 0$.

\subsubsection{Hypergeometric functions.}
Let $r\le s+1$.
Let $a_1,\dots,a_r\in\C$ and $b_1,\dots,b_s\in\C\setminus\Z_-$.
Generalized hypergeometric functions are defined by
\begin{equation}
\hyp{r}{s}\left[\genfrac{}{}{0pt}{}{a_1,\dots,a_r}{b_1,\dots,b_r};z\right]
:=\sum_{m=0}^\infty
\frac{(a_1)_m\cdots (a_r)_m}{m! (b_1)_m\cdots (b_s)_m}\cdot z^n,
\label{eq:rFs}
\end{equation}
where $(a)_m:=a(a+1)\cdots(a+m-1)$ is the Pochhammer symbol.
For $r\le s$ the radius of convergence is $\infty$, for $r=s+1$ it is $1$.
For $r>s+1$ the series diverges.

\subsubsection{Notation for lists}
Let $a_1, \dots, a_p\in \C$ and $h\in\C$.
We use the following notation for lists,
\begin{gather*}
(a):= a_1, \dots, a_p;\\
(a)+h:= a_1+h, \dots, a_p+h;\\
(a)_{\setminus j}:= a_1,\dots,a_{j-1},a_{j+1},\dots,a_p.
\end{gather*}
In particular, we denote \eqref{eq:rFs}
by $\hyp{r}{s}\big[\genfrac{}{}{0pt}{}{(a)}{(b)};z\big]$.

\subsubsection{Double powers}
Denote by $\C^\times$ (resp.\ $\R^\times_+$) the multiplicative group of $\C$
(resp.\ the multiplicative group of positive reals).
By $\T\subset \C^\times$ we denote the subgroup
$|z|=1$, we have $\T\simeq\R/2\pi \Z$,
\begin{equation*}
\C^\times= \T\times \R^\times_+,
%\label{eq:times1}
\end{equation*}
and $\R^\times_+$ is isomorphic to the additive group $\R$ of real numbers.

The dual group $(\C^\times)^*$, i.e., the group of all homomorphisms
$\C^\times\to \T$, is
\begin{equation*}
(\C^\times)^*= (\T)^*\times (\R^\times_+)^*\simeq \Z\times \R.
%\label{eq:times2}
\end{equation*}

Denote by $\Lambda_\C$ the set of pairs
$a\big| a'$ such that $a$, $a'\in\C$ and $a-a'\in\Z$.
We write such pairs by
\[
a\big|a'=\tfrac{k+\sigma}2\big|\tfrac{-k+\sigma}2,
\qquad\text{where $k\in\Z$, $\sigma\in\C$}.
\]
By $\Lambda\subset \Lambda_\C$ we denote the subset
consisting of $a|a'$ satisfying $a'+\ov a=0$,
\[
a\big|a'=
\tfrac{k+{\rm i}s}2\big|\tfrac{-k+{\rm i}s}2\in \Lambda,\qquad\text{where $s\in \R$}.
\]

Define the following functions on $\C^\times$:
\[
z^{a|a'}:= z^a \ov z^{\,a'}:= (z/\ov z)^{k/2} |z|^{\sigma}.
\]
Notice that such functions are precisely all homomorphisms
$\C^\times\to \C^\times$.
For $a|a'\in \Lambda$ we get homomorphisms $\C^\times\to \T$,
\[
\Big|z^{\tfrac{k+{\rm i}s}2\big|\tfrac{-k+{\rm i}s}2} \Big|=1.
\]

\subsubsection{The complex Mellin transform}
Denote by $\dd t$ the Lebesgue measure on $\C$,
\[
\dd t:= {\rm d}\Re t\, {\rm d}\Im t.
\]

We define the complex Mellin transform $\cM$ as the Fourier transform on the group $\C^\times$.
For a~function $f$ on $\C^\times$ we define the Mellin transform as a function
on $\Lambda_\C$ defined by
\[
g\big(a\big|a'\big)=
\cM f\big(a\big|a'\big):=\frac 1{2\pi}
\int_{\C^\times} t^{a|a'} f(t)\frac {\dd t}{t^{1|1}},
\]
the factor ${\dd t}/{t^{1|1}}$ is the $\C^\times$-invariant measure on $\C^\times$.

The Mellin transform is a unitary operator from
$L^2\big(\C^\times,{\dd t}/{t^{1|1}} \big)$
to $L^2$ on $\Lambda\simeq \Z\times\R$, the inversion formula is
\begin{equation*}
f(t)=\frac 1{2\pi}\sum_{k\in\Z} \int_\R g\big(\tfrac{k+{\rm i}s}2\big|\tfrac{-k+{\rm i}s}2\big) t^{{\tfrac{-k-{\rm i}s}2\big|\tfrac{k-{\rm i}s}2} }\,{\rm d}s.
%\label{eq:inversion}
\end{equation*}

The convolution on the group
$\C^\times$ is defined by formula
\[ f*g(z):=\int_\C f(t)g(z/t)\,\frac{\dd t}{t^{1|1}}.
\]
The Mellin transform sends convolutions to products,
\[
\cM (f*g)(z)=2\pi\, \cM f(z)\cdot \cM g(z).
\]

\subsection{The Gamma-function of the complex field}\label{ss:Gamma}
Following Gelfand, Graev and Retakh \cite{GGR},
we define the Gamma-function of the complex field as
a function on $\Lambda_\C$ by\footnote{The definition slightly
differs from a definition from \cite{GGR}, which was used in
\cite{MN,Ner-doug}. We write $\Im t$ instead of $\Re t$. For this reason
the factor $i^{a-a'}$ from \cite{GGR,MN} disappears.}
\begin{align}
\Gamma^\C\big(a\big|a'\big)&:=
\frac 1\pi \int_\C t^{a-1|a'-1} {\rm e}^{2{\rm i} \Im t}\dd ta \notag\\
&\hphantom{:}=
\frac{1}{\pi} \lim_{r\to\infty} \int_{|t|\le r} t^{a-1|a'-1}
{\rm e}^{2{\rm i} \Im t}\dd t
= \frac{\Gamma(a)}{\Gamma(1-a')}
=\frac{(-1)^{a-a'}\Gamma(a')}{\Gamma(1-a)} \notag \\
&\hphantom{:}= \frac{1}{\pi}\,
\Gamma(a)\Gamma(a')\sin\pi a'
= \frac{(-1)^{a-a'}}{\pi} \,\Gamma(a)\Gamma(a')\sin\pi a. \label{eq:Gamma}
\end{align}
The integral conditionally converges if $0<\Re(a+a')<1$
and diverges otherwise. Clearly, the right hand side is meromorphic
in the whole $\Lambda_\C$.

\begin{Remark} In particular, we can write \eqref{eq:beta} as
\[
\B^\C(\alpha|\alpha',\beta|\beta')=
\frac{\Gamma^\C(\alpha|\alpha')\,\Gamma^\C(\beta|\beta')}
{\Gamma^\C(\alpha+\beta|\alpha'+\beta')}.
\]
\end{Remark}

It is easy to see that
\begin{subequations}
\begin{gather}
\Gamma^\C\big(a|a'\big)=(-1)^{a-a} \Gamma^\C\big(a'|a\big),\label{eq:symmetry} \\
\Gamma^\C\big(a|a'\big)\, \Gamma^\C\big(1-a|1-a'\big) =(-1)^{a-a'},\label{eq:gamma-complement} \\
\Gamma^\C\big(a+m\big|a'+m'\big) =(-1)^{m'} \Gamma^\C\big(a\big|a'\big)(a)_m(a')_{m'}, \label{eq:shift1} \\
 \Gamma^\C\big(a-m\big|a'-m'\big) =\frac{(-1)^{m} \Gamma^\C\big(a\big|a'\big)}{(1-a)_m(1-a')_{m'}}, \label{eq:shift2} \\
\prod_{j=0}^{k-1}\Gamma^\C\big(a+\tfrac jk\big|a'+\tfrac jk \big)=\Gamma^\C\big(ka\big|ka'\big)k^{1-(a+a')k},\label{eq:gamma-gauss}
\end{gather}
\end{subequations}
where $m\in\N$, $k\in\N$.
Values of $\Gamma^\C$ at integer points are
\begin{gather*}
 \Gamma^\C(k|l) = 0,\qquad \text{for $k$, $l\in\N$;} \\
 \Gamma(m|-k) =\frac{(m-1)!(-1)^k}{k!},\qquad \text{for $m\in \N$, $k\in \Z_+$}.
\end{gather*}
For $m$, $m'\in\Z_+$ we have a pole at $(-m)|(-m')$, more precisely,
\begin{equation}
\res\limits_{\epsilon=0}
\Gamma^\C(-m+\epsilon|-m'+\epsilon)=\frac{(-1)^m}{m!\, m'!}.
\label{eq:res-gamma}
\end{equation}
Next (see Section~\ref{ss:asy} below),
for $a|a'\in\Lambda_\C$ and $\xi=\frac12(k+{\rm i}s)$
with $k\in\Z$, $s\in \R$ we have
\begin{equation}
\Gamma^\C\big(a+\xi|a'-\ov \xi\big)=
\exp\big\{2{\rm i}\Im (\xi\ln \xi-\xi)\big\}\cdot \xi^{a-\tfrac12\big|a'-\tfrac12} \big(1+O\big(|\xi|^{-1}\big)\big).
\label{eq:asy1}
\end{equation}
Therefore,
\begin{equation}
\big| \Gamma^\C\big(a+\xi|a'-\ov \xi\big)\big|\sim |\xi|^{\Re(a+a')-1}.
\label{eq:asy2}
\end{equation}

\subsection{Hypergeometric functions of the complex field}\label{ss:def}
Let $a_1|a'_1,\dots, a_q|a'_q$, $b_1|b'_1,\dots,b_p|b'_p\in \Lambda_\C$.
Temporarily, we assume that they satisfy the conditions:
\begin{equation}
\Re(a_\alpha+a_\alpha')>0,\qquad \Re(b_\beta+b_\beta')>0.
\label{eq:cond1}
\end{equation}
and
\begin{equation}
\sum_\alpha \Re(a_\alpha+a_\alpha')+\sum_\beta\Re(b_\beta+b_\beta')<p+q.
 \label{eq:cond2}
\end{equation}

We define the following function on $\Lambda_\C$:
\begin{gather}
\Kc{p}{q}\left[\genfrac{}{}{0pt}{}{(a|a')}{(b|b')};  \tfrac{k+\sigma}{2}\Bigl| \tfrac{-k+\sigma}{2}
\right]\nonumber\\
\qquad{} :=
 \prod_{\alpha=1}^q \Gamma^\C\Bigl(a_\alpha+\tfrac{k+\sigma}2\Bigl|a_\alpha'+\tfrac{-k+\sigma}2\Bigr)
 \prod_{\beta=1}^p \Gamma^\C\Bigl(b_\beta+\tfrac{-k-\sigma}2\Bigl|b_\beta'+\tfrac{k-\sigma}2\Bigr).
\label{eq:cK}
\end{gather}
Next, we define the hypergeometric functions of the complex field
as the contour integral:
\begin{equation}
\Gc{p}{q}\left[\genfrac{}{}{0pt}{}{(a|a')}{(b|b')};z\right]
:=\frac{1}{2\pi {\rm i}}
\sum_{k\in\Z} \int_{{\rm i}\R}
\Kc{p}{q}\left[\genfrac{}{}{0pt}{}{(a|a')}{(b|b')};\tfrac{k+\sigma}{2}\Bigl| \tfrac{-k+\sigma}{2}\right]
 z^{\frac{-k-\sigma}2\big| \frac{k-\sigma}2} \,{\rm d}\sigma.
 \label{eq:def}
\end{equation}
The integration is taken along the imaginary axis ${\rm i}\R$.
The condition \eqref{eq:cond2} provides the conditional convergence of
the integral in $\sigma$ and the absolute convergence of the series
(see Sections~\ref{ss:meijer} and \ref{ss:evaluation} below).
Under the stronger condition
\begin{equation}
\sum_\alpha \Re(a_\alpha+a_\alpha')+\sum_\beta\Re(b_\beta+b_\beta')<p+q-1
 \label{eq:cond3}
\end{equation}
all integrals in $\sigma$ convergence absolutely
(this follows from \eqref{eq:asy2}).

The expression \eqref{eq:cK} has poles (they are analyzed in
Section~\ref{ss:evaluation}) in $\sigma$ originating from two groups
of factors.
We call poles of the factors
$\Gamma^\C\big(a+\frac{k+\sigma}2\big|a'+\frac{-k+\sigma}2\big)$
\emph{left poles}, the poles of the factors
$\Gamma^\C\big(b+\frac{-k-\sigma}2\big|b'+\frac{k-\sigma}2\big)$
\emph{right poles}.
Under the conditions \eqref{eq:cond1} left poles are contained
in the left half-plane $\Re\sigma<0$,
right poles are contained in the right half plane $\Re\sigma>0$, and
the axis ${\rm i}\R$ separates these groups of poles.
